% 数论（综述）
% license CCBY4
% type Wiki

本文根据 CC-BY-SA 协议转载翻译自维基百科\href{https://en.wikipedia.org/wiki/Number_theory}{相关文章}

\begin{figure}[ht]
\centering
\includegraphics[width=6cm]{./figures/ae528936f6668ed9.png}
\caption{} \label{fig_SHuLun_1}
\end{figure}
数论是纯数学的一个分支，主要致力于研究整数以及算术函数。数论学家研究素数，以及由整数构造的数学对象（例如有理数）的性质，或作为整数的推广而定义的对象（例如代数整数）。

整数既可以单独研究，也可以作为某些方程的解来研究（丢番图几何）。数论中的问题常常能够通过研究某些解析对象来理解，例如黎曼ζ函数，它以某种方式编码了整数、素数或其他数论对象的性质（解析数论）。另外，还可以研究实数与有理数之间的关系，例如研究无理数如何用分数来逼近（丢番图逼近）。

数论与几何一样，是数学最古老的分支之一。数论的一个特点是，它处理的往往是非常容易理解却极难解决的命题。例如，费马大定理在最初提出后的 358 年才被证明，而哥德巴赫猜想自 18 世纪提出以来至今仍未解决。德国数学家卡尔·弗里德里希·高斯曾经说过：“数学是科学的女王，而数论是数学的女王。”\(^\text{[1]}\)数论一度被视为纯粹数学的典范，认为它在数学之外没有任何应用，直到 20 世纪 70 年代，人们才发现素数可被用作构建公钥密码算法的基础。
\subsection{历史}

数论是数学的一个分支，研究整数及其性质与关系。[2] 整数是一个集合，它在自然数集合${1,2,3,\dots}$的基础上扩展，加入了数0，以及自然数的相反数${-1,-2,-3,\dots}$。数论学家不仅研究素数，还研究由整数构造的数学对象（例如有理数）的性质，或定义为整数推广的对象（例如代数整数）。[3][4]

数论与算术密切相关，一些作者甚至将这两个词视为同义词。[5] 然而，今天“算术”一词通常指研究数值运算的学科，并扩展到实数范围。[6] 在更具体的意义上，数论仅限于研究整数，关注它们的性质和相互关系。[7] 传统上，它也被称为高等算术。[8] 到 20 世纪初，“数论”这一术语已被广泛采用。[note 1]其中“number”意为整数，既可以指自然数，也可以指全体整数。[9][10][11]

**初等数论**研究整数的一些方面，这些方面可以用初等方法来探讨，例如利用初等证明。[12] 相比之下，**解析数论**依赖于复数以及来自分析学和微积分的技巧。[13] **代数数论**则使用代数结构（如域和环）来研究数的性质及其相互关系。**几何数论**运用几何学的概念研究数。[14] 数论的其他分支还包括：**概率数论**[15]、**组合数论**[16]、**计算数论**[17]，以及**应用数论**，后者研究数论在科学与技术中的应用。[18]
